% paper/intro_related.tex
% Skeletons for the Introduction and Related Work sections.
% Self-contained: paste straight into your paper template.

\section{Introduction}\label{sec:intro}

Large language models (LLMs) are increasingly used as the cognitive
core of embodied agents — they read a goal in natural language, look
at a list of available actions, and emit a plan that a low-level
controller executes~\cite{ahn2022saycan, huang2022innermonologue,
singh2022progprompt, liang2022codeaspolicies}. The headline result is
striking: even modest open-weight LLMs achieve high \emph{task
success} on benchmarks like ALFWorld~\cite{shridhar2020alfworld} and
AI2-THOR~\cite{kolve2017thor}. Closer inspection, however, reveals a
hidden cost. Plans produced by current methods average $2$--$3\times$
more steps than the ground-truth reference\footnote{Measured on
pyplanner's 38-task AI2-THOR set;
see \texttt{record\_report.json} in the artifact.}, and execution
``succeeds'' only because most agents include a generous retry-on-failure
loop that masks brittle plans behind extra LLM calls.

We argue that two design choices are responsible. First, planners do
not verify plans against a grounded model of the environment before
execution — preconditions like \emph{``you must Find an object before
Picking it''} are encoded only in the LLM's prompt and enforced only
when the simulator returns an error. Second, on failure, planners
typically replan the entire remaining task, paying the full token cost
of re-explaining context that has already been correctly handled. Both
choices are tractable to fix without retraining the LLM.

We propose \textbf{GRACE} — Grounded Repair-and-Critique Embodied
planner — a single-LLM planner that integrates three components:
\begin{enumerate}
  \item a \textbf{rule-based symbolic precondition verifier}
        that runs in $O(|\pi|)$ time and zero tokens,
  \item \textbf{hierarchical decomposition with sub-goal--localised
        repair} — when a step fails, only the suffix of the failed
        sub-goal is regenerated, and
  \item a \textbf{hybrid memory store} that combines a small curated
        seed pool of ground-truth plans with a growing pool of
        successful runtime episodes.
\end{enumerate}

Our contributions are:

\begin{itemize}
  \item GRACE, an integrated planner whose three components are
    each lightweight enough to deploy on a single Ollama server
    (\S\ref{sec:method}).
  \item A symbolic precondition verifier written in
    $\sim 250$ lines of Python that operates on the canonical
    pyplanner action vocabulary and emits typed violation reasons
    suitable for direct prompt feedback (\S\ref{sec:verifier}).
  \item An evaluation against eight published planning strategies on
    the 38-task AI2-THOR benchmark across four model sizes,
    plus a three-way ablation that isolates the contribution of each
    component (\S\ref{sec:experiments}).
\end{itemize}

\section{Related Work}\label{sec:related}

\paragraph{LLM-as-planner.}
SayCan~\cite{ahn2022saycan} grounds skill selection in a learned
affordance value; Inner Monologue~\cite{huang2022innermonologue} adds
textual feedback after each step.
ProgPrompt~\cite{singh2022progprompt} and Code-as-Policies~\cite{liang2022codeaspolicies}
push the LLM toward program-style output. LLM+P~\cite{liu2023llmp}
translates NL to PDDL and delegates to a classical planner;
LLM-DP~\cite{dagan2023llmdp} maintains a symbolic world model
alongside the LLM. GRACE differs by keeping the planner LLM-driven
but inserting a deterministic check between plan and execution.

\paragraph{Reasoning, refinement, and self-critique.}
Chain-of-Thought~\cite{wei2022cot} elicits intermediate reasoning;
ReAct~\cite{yao2022react} interleaves thought and action;
Tree-of-Thoughts~\cite{yao2023tot} searches over reasoning branches.
Self-Refine~\cite{madaan2023selfrefine} and Reflexion~\cite{shinn2023reflexion}
critique with the same LLM. GRACE replaces the LLM critic with a
zero-cost symbolic checker that emits typed feedback — sidestepping
the well-known LLM self-evaluation bias.

\paragraph{Embodied planning with verification.}
LLM-Planner~\cite{song2023llmplanner} retrieves few-shot examples by
embedding similarity. AdaPlanner~\cite{sun2023adaplanner} distinguishes
in-plan refinement from out-of-plan replanning.
ISR-LLM~\cite{zhou2024isrllm} validates LLM-emitted PDDL with a
classical solver. The LLM-Modulo framework~\cite{kambhampati2024llmmodulo}
generalises ``LLM-generate + external-verifier'' to a wide range of
tasks. GRACE specialises this idea to embodied household tasks with a
\emph{hand-written, vocabulary-specific} verifier that costs orders of
magnitude less than a classical-planner verifier.

\paragraph{Memory-augmented agents.}
Voyager~\cite{wang2023voyager} stores reusable code skills;
ExpeL~\cite{zhao2024expel} distills natural-language insights from
past trajectories; AutoGuide~\cite{fu2024autoguide} produces
state-conditioned guidelines. GRACE's memory differs by being
\emph{hybrid}: a curated seed pool gives cold-start coverage, while a
live pool accumulates successful plans over time, both surfaced via
the same prompt-block format.

\paragraph{Hierarchical and repair-oriented planning.}
HiP~\cite{ajay2023hip} and the pyplanner Hierarchical and
HierarchicalFewShot strategies decompose tasks into sub-goals.
DEPS~\cite{wang2023deps} introduces failure description + sub-goal
selection in Minecraft. ReplanVLM~\cite{mei2024replanvlm} uses a VLM
to detect failures and trigger replanning. To our knowledge GRACE is
the first to (a) retain sub-goal boundaries \emph{at execution time}
and (b) restrict repair to the suffix of the failed sub-goal, so that
already-completed work and untouched future sub-goals are preserved.
